\begin{abstract}
Home-Cage Monitoring (HCM) systems observe laboratory animals continuously in their own cage, improving the contextual richness and reproducibility of preclinical data. Yet HCM systems are highly heterogeneous --- differing in sensors, measured variables, software, and output formats --- and their data are inseparable from a rich experimental context. This heterogeneity hinders comparison, reuse, and FAIR data sharing. We present HCMO, the Home-Cage Monitoring Ontology, which builds on earlier conceptual models of HCM and, to our knowledge, is the first machine-actionable ontology to represent HCM as an integrated experimental and data-acquisition system. HCMO is organised into five modules covering monitored enclosures, biological subjects, environmental context, observations and results, and technical acquisition. It is built around an explicit separation of sensor, observation, and result, preserving the chain from device to interpreted measurement. The published HCMO 0.3.0 release reuses BFO, IAO, the 2017 SOSA/SSN Recommendation, Schema.org, OWL-Time, and a reviewed SemTS 1.2.0 time-series segment pattern where their semantics fit. Provenance-backed state records and a pinned QUDT 3.4.0 quantity pattern are implemented; broader domain alignments remain future work. HCMO is delivered as a FAIR, tool-consumable package: modular sources, reproducibly generated distributions (Turtle/OWL/JSON-LD), SHACL shapes, competency-question queries, a JSON-LD context, persistent identifiers, an open licence, and documentation. We describe its design, evaluate it through competency questions and ontology-quality tooling, and discuss its uptake within the COST TEATIME community.

\keywords{ontology \and home-cage monitoring \and laboratory animals \and FAIR data \and SOSA/SSN \and 3Rs}
\end{abstract}
